\section{Members and Activities}\label{sec:members}

The group consists of Lalia Sadio, Adriana Souza, Leandro Bergmann and Dalila Rezende. Table~\ref{tab:members} lists the activities of each member.

\begin{table}[h]
\centering
\caption{Activities carried out by each group member.}\label{tab:members}
\small
\begin{tabularx}{\linewidth}{@{}l X@{}}
\toprule
Member & Activities\\
\midrule
Lalia Sadio & Proposed the topic of the project. Algorithms, together with Leandro: implementation of the algorithms (\code{algorithms/}) and of the experimental code, and the sections of this report that describe, analyze and test them (Sections~\ref{sec:design}--\ref{sec:experiments}).\\
Leandro Bergmann & Algorithms, together with Lalia: implementation of the algorithms (\code{algorithms/}) and of the experimental code, and the sections of this report that describe, analyze and test them (Sections~\ref{sec:design}--\ref{sec:experiments}).\\
Adriana Souza & Theoretical background, together with Dalila: the problem, the techniques and the review of the recent related work (Section~\ref{sec:background}).\\
Dalila Rezende & Theoretical background, together with Adriana, having the Sebrae learning platform (GuIA) as its basis (Section~\ref{sec:background}).\\
All members & Every other activity of the project, carried out and reviewed together: real-data experiments, GRASP, independent validation, discussion, conclusion, and the presentation. Each member reviewed the whole code and the whole report.\\
\bottomrule
\end{tabularx}
\end{table}

All members reviewed the whole code and the whole report, and each one is able to explain every code segment and every complexity derivation in the presentation, as required by the assignment.

\paragraph{Use of AI tools.} Claude (Anthropic) was used as an assistant for the English review of this report and for reviewing the code: an edge-case fix in \code{grasp.py} (\code{alpha}\,$=0$), comments and docstrings, formatting, and the unit test of the worked example. It also helped to prepare drafts of the presentation material. The group reviewed every change and checked it by re-running the experiments, which produced identical values and operation counts before and after the review.

\section{Conclusion}\label{sec:conclusion}

\paragraph{General comments.} We studied the selection of learning modules under a study-time budget, modeled as the 0/1 knapsack problem, starting from the article RUTICO, whose solution is an exhaustive DFS. We showed that dynamic programming is the same search with memoization: the decision tree has $2^n$ leaves but only $n(T+1)$ distinct states. We proved the correctness and the best-, worst- and space complexity of the DFS ($\Theta(n)$ to $\Theta(2^n)$ nodes, with exactly $2^{n+1}-1$ nodes when everything fits), of the DP ($\Theta(nT)$, pseudo-polynomial) and of the greedy ($\Theta(n\log n)$, without a constant guarantee unless the best single module is considered), and of three additional algorithms. The experiments, based on hardware-independent operation counters and validated on more than a thousand random instances, confirmed each of these results. The DP is the most suitable algorithm for the learning-path scenario because it is exact, has a cost that does not depend on the data, and the budget $T$ is small in practice; the experiments also showed regimes in which Branch and Bound, Meet-in-the-Middle or the greedy are preferable. The comparison was then taken to anonymized real data from a platform in production (12 modules and 78 topics, Section~\ref{sec:real}), which confirmed the correlation hypothesis for the content score ($r=+0.92$), showed a negatively correlated regime for the completion score ($r=-0.62$) that the synthetic families do not cover, and showed the greedy losing up to 21\% of the score with short budgets. The complementary experiment with GRASP (Section~\ref{sec:grasp}) covered the missing metaheuristic technique and, by costing more and losing more than the DP on the topic instance, reached in e-learning a conclusion aligned with that of MAS-Cloud+ in the cloud: when the problem fits the exact method, it is the best choice. An independent validation against OR-Tools and SciPy/HiGHS on 3,533 instances found no divergences (Section~\ref{sec:validation}). Compared with recent learning-path work (Section~\ref{sec:related}), the contribution is not a new recommender but a measured answer to what the selection step costs and which technique to use.

\paragraph{Main difficulties.}
\begin{itemize}[leftmargin=1.4em,itemsep=2pt]
\item \textbf{The article does not use dynamic programming.} The solution of the authors is a DFS, so we had to implement it and justify the DP through its relation to the DFS (memoization) rather than assume it.
\item \textbf{Bridging the article and the knapsack.} The article works on a graph with prerequisites; the knapsack ignores them. Deciding how to state this relaxation, and keeping a graph version (\code{graph\_dfs.py}) for the authors' setting, required care.
\item \textbf{Fair comparisons.} Comparing algorithms with different cost structures is delicate: a Branch and Bound node costs $\bigO(n)$ because of the bound, a DP cell costs $\bigO(1)$. We counted the bound steps explicitly and we did not rely on wall-clock time alone.
\item \textbf{Interpreting the greedy.} Its mean gap is small for large $n$ but its worst gap is large; reporting only the mean would have been misleading, and showing only the worst case would have been unfair. The families and the budget ratio were needed to tell the whole story.
\item \textbf{Exponential experiments.} The DFS cannot be run beyond about 22 modules in reasonable time, so growth beyond that point is established by proof and extrapolation, not by measurement.
\item \textbf{Real data and their proxies.} The platform measures neither learning nor the value of a module, so the scores are proxies and the topic times come from a declared proportional split; extracting the data required read-only aggregated queries and full anonymization, and the instances are small.
\item \textbf{An independent check revealed what the experiments did not.} Only the validation with external libraries, at larger sizes, exposed that Branch and Bound may not finish on strongly correlated instances (Section~\ref{sec:validation}), a behavior invisible up to $n=40$.
\item \textbf{Implementation details.} Hand-writing the merge sort to keep $\Theta(n\log n)$ in every case, using integer arithmetic for ratio comparisons, and handling deep recursion in the memoized DP.
\end{itemize}

\paragraph{Limitations.} The relaxation without prerequisites, synthetic instances and small real instances from a single platform, proxy scores, a single machine and interpreter, a targeted (not systematic) review of the related work, and instance sizes up to $n=600$ (and $n=40$ for the exponential algorithms) bound the scope of the conclusions (Section~\ref{sec:discussion}).

\paragraph{Future work.} (i) Incorporate prerequisites: the precedence-constrained knapsack, starting from tree-shaped prerequisites, for which dynamic programming remains applicable. (ii) Several constraints at once (time and difficulty), i.e.\ a multidimensional knapsack, where the DP table acquires one dimension per constraint. (iii) Replace the proxy scores by the time and score estimates of the article \cite{nabizadeh2017rutico,nabizadeh2019uncertainty} or by the output of learner models such as \cite{lv2025genal,yang2026forgetting}, and extend the real-data study to more platforms. Integer-programming formulations such as that of \cite{bose2026learnopt} are a natural baseline for the version with prerequisites. (iv) Evaluate larger instances with a compiled implementation, and mixed strategies such as DP for the bulk of the budget and Branch and Bound for the remainder.
